% Original Markdown SHA-256: 45c4e09b089b65afe65ae661494c4b82bf41e8f95b76971421f8243eec6f3894
\chapter{General-k exponential capacity, exact carry certificates, and improved k=5,6 constructions}
\label{chapter:02_general-k-capacity}
{\small\noindent\textbf{Source:} \nolinkurl{history/EP817_CUMULATIVE_PR2_PR3_20260914/PR3_general_k/notes/general-k-capacity.md}\
\textbf{Identity:} \nolinkurl{45c4e09b089b65afe65ae661494c4b82bf41e8f95b76971421f8243eec6f3894}}
\medskip
14 September 2026. Research continuation in the Erdős Problem 817 workbench. Base revision: \nolinkurl{23c0110c95b5a2036bdc04a1f352b6e5e27742c8}.

\textbf{Evidence status.} The general statements below have ordinary mathematical proofs. The finite blocks have independently replayed integer certificates. These claims are submitted for independent review; no new Lean elaboration is claimed. Existing Core/Extended sources, their receipts, PR \#1, and the separate variance PR \#2 remain unchanged.

\section{1. Problem, attribution, and results}\label{n02-problem-attribution-and-results}

For a finite set A of distinct positive integers, write

\[
\begin{adjustbox}{max width=.92\linewidth}
$\displaystyle H(A)=\left\{\sum_{a\in U}a:U\subseteq A\right\},\qquad S(A)=\sum_{a\in A}a.$
\end{adjustbox}
\]

The empty sum belongs to H(A). For k\textgreater=3, call A k-admissible when H(A) contains no tuple (x,x+d,...,x+(k-1)d) with d a nonzero integer. Define

\[
\begin{adjustbox}{max width=.92\linewidth}
$\displaystyle g_k(n)=\min\{\max A:|A|=n,\ A\text{ is k-admissible}\}\quad(n\ge1).$
\end{adjustbox}
\]

These are the conventions of Erdős--Sárközy and Korsky. The anonymous/deleted Reddit contributor supplied the workbench\textquotesingle s original k=4 construction and signed-block proof. The Lean finite-upper-bound extension is separately credited to sneed-and-feed. The present general-k arguments do not assume the proposed variance refinement in PR \#2; the original k=4 block is used as a control example. No claim of literature priority is made for the arguments here.

The continuation establishes the following.

\textbf{Theorem 1 (all-k capacity).} For every fixed integer k\textgreater=3, the limit

\[
\begin{adjustbox}{max width=.92\linewidth}
$\displaystyle \Lambda_k=\lim_{n\to\infty}g_k(n)^{1/n}$
\end{adjustbox}
\]

exists. It has the exact finite-certificate characterization

\[
\begin{adjustbox}{max width=.92\linewidth}
$\displaystyle \boxed{\Lambda_k=
 \inf_{\substack{B\ne\varnothing\text{ finite, distinct, positive}\\
 q>\sum B;\ H(B)\bmod q\text{ has no nonzero-step k-AP}}}
 q^{1/|B|}.}$
\end{adjustbox}
\]

The modulus q ranges over all integers, including composite integers. A modular progression here is an ordered tuple with nonzero step in Z/qZ; repetition of residues is allowed. Every test in the displayed infimum is finite and decidable.

\textbf{Theorem 2 (all-length decision).} For any finite set D of nonnegative integers containing 0, any q\textgreater=2, and any k\textgreater=3, a finite directed graph decides whether

\[
\begin{adjustbox}{max width=.92\linewidth}
$\displaystyle L_t(D,q)=\left\{\sum_{j=0}^{t-1}d_jq^j:d_j\in D\right\}$
\end{adjustbox}
\]

is k-progression-free for every t. Digits \textgreater=q and signed carries are retained. An unsuccessful certificate returns an explicit integer progression and its digit representations. When D is contained in {[}0,q-1{]}, the graph has at most 2k! states.

\textbf{Theorem 3 (explicit improved bounds).} Put W=(1,4,5,17,21,22). For n\textgreater=1 write

\[
\begin{adjustbox}{max width=.92\linewidth}
$\displaystyle j=\lfloor(n-1)/6\rfloor,\qquad s=1+((n-1)\bmod6).$
\end{adjustbox}
\]

With w\_s the s-th entry of W,

\[
\begin{adjustbox}{max width=.92\linewidth}
$\displaystyle \boxed{g_5(n)\le w_s97^j,\qquad g_6(n)\le w_s93^j.}$
\end{adjustbox}
\]

In particular,

\[
\begin{adjustbox}{max width=.92\linewidth}
$\displaystyle \Lambda_5\le97^{1/6},\qquad \Lambda_6\le93^{1/6}.$
\end{adjustbox}
\]

These improve the corresponding upper rates in Korsky\textquotesingle s inspected v1, namely sqrt(5) and 7\^{}(2/5). The exact comparisons are 97\textless5\^{}3 and \(93^5<7^{12}\). No floating-point comparison is needed.

Theorem 1 identifies the full fixed-k exponential-rate problem with a uniform bound on finite certificate capacities. Determining that infimum numerically for every k, or establishing the sharp large-k law, remains a separate mathematical task. Section 10 gives the exact outstanding inequality rather than an assumed premise for any theorem above.

\section{2. The four-term obstruction and its exact witness map}\label{n02-the-four-term-obstruction-and-its-exact-witness-map}

For each k\textgreater=4 there is an inclusion of sets of admissible generator sets

\[
\begin{adjustbox}{max width=.92\linewidth}
$\displaystyle \mathcal A_{4,n}\hookrightarrow\mathcal A_{k,n},\qquad A\longmapsto A.$
\end{adjustbox}
\]

A nonconstant k-AP would contain a nonconstant 4-AP in its first four positions. This proves the inclusion and transfers upper constructions in that direction.

The inclusion is strict already for k=5 and A=\{1,2\}. Its subset-sum image is \{0,1,2,3\}, so it is 5-admissible. The short relation 2*1-2=0 has nonzero magnitude-one and magnitude-two layers. Feeding its literal coefficient classes into the four-term witness construction gives

\[
\begin{adjustbox}{max width=.92\linewidth}
$\displaystyle P_1=0,\quad P_2=1,\quad N_1=2,\quad N_2=0,$
\end{adjustbox}
\]

and the tuple

\[
\begin{adjustbox}{max width=.92\linewidth}
$\displaystyle (P_1+N_2,P_1+P_2,N_1+N_2,N_1+P_2)=(0,1,2,3).$
\end{adjustbox}
\]

Thus the map from this coefficient relation to a four-term progression works exactly as before, but its codomain is the allowed four-term witness set for a 5-admissible object. The general-k development below retains simultaneous progression equations instead of asserting that these coefficient layers vanish.

\section{3. An exact composition of admissible generator sets}\label{n02-an-exact-composition-of-admissible-generator-sets}

Let A and B be finite sets of distinct positive integers, and put

\[
\begin{adjustbox}{max width=.92\linewidth}
$\displaystyle R=2S(A)+1,\qquad C=A\cup RB.$
\end{adjustbox}
\]

All members of RB exceed every member of A, so \textbar C\textbar=\textbar A\textbar+\textbar B\textbar. At the level of chosen subsets, the map

\[
\begin{adjustbox}{max width=.92\linewidth}
$\displaystyle \mathcal P(A)\times\mathcal P(B)\longrightarrow\mathcal P(C),
 \qquad(U,V)\longmapsto U\cup RV$
\end{adjustbox}
\]

is a bijection. Its inverse intersects with A and RB and divides the latter elements by R. It induces the following value-space bijection:

\[
\begin{adjustbox}{max width=.92\linewidth}
$\displaystyle E_{A,B}:H(A)\times H(B)\longrightarrow H(C),\qquad(x,y)\longmapsto x+Ry.$
\end{adjustbox}
\]

Indeed, 0\textless=x\textless=S(A)\textless R, so the inverse on H(C) is

\[
\begin{adjustbox}{max width=.92\linewidth}
$\displaystyle z\longmapsto(z\bmod R,\lfloor z/R\rfloor).$
\end{adjustbox}
\]

Collisions among different subsets within H(A) or H(B) remain present; this statement is about the distinct numerical image sets.

\textbf{Lemma 3.1 (progression decomposition).} Every k-term progression in H(C) pulls back under E to a k-term progression in each factor.

\textbf{Proof.} Write its terms u\_i=x\_i+Ry\_i. For 0\textless=i\textless=k-3,

\[
\begin{adjustbox}{max width=.92\linewidth}
$\displaystyle 0=u_i-2u_{i+1}+u_{i+2}
 =(x_i-2x_{i+1}+x_{i+2})+R(y_i-2y_{i+1}+y_{i+2}).$
\end{adjustbox}
\]

The first parenthesis is an integer in {[}-2S(A),2S(A){]}={[}-(R-1),R-1{]}. It is divisible by R and hence equals zero. The second parenthesis therefore also equals zero. These consecutive second-difference identities prove that both coordinate sequences are arithmetic progressions. QED.

The corresponding constructive contrapositive is a map

\[
\begin{adjustbox}{max width=.92\linewidth}
$\displaystyle \operatorname{AP}_k^*(H(C))\longrightarrow
 \operatorname{AP}_k^*(H(A))\sqcup\operatorname{AP}_k^*(H(B)).$
\end{adjustbox}
\]

Pull back the tuple; send it to the first component when its first common difference is nonzero, otherwise to the second. The second difference is then nonzero because the original difference is d\_x+Rd\_y. Thus every failure of the composition has an explicit failure witness in a factor.

\textbf{Corollary 3.2.} If A and B are k-admissible, then C is k-admissible. Furthermore,

\[
\begin{adjustbox}{max width=.92\linewidth}
$\displaystyle 2S(C)+1=(2S(A)+1)(2S(B)+1).$
\end{adjustbox}
\]

The equality follows by substituting S(C)=S(A)+(2S(A)+1)S(B); every term and constant is retained.

\section{4. Proof that every fixed-k exponential rate exists}\label{n02-proof-that-every-fixed-k-exponential-rate-exists}

Define

\[
\begin{adjustbox}{max width=.92\linewidth}
$\displaystyle F_k(n)=\min_{A\in\mathcal A_{k,n}}(2S(A)+1),\qquad F_k(0)=1.$
\end{adjustbox}
\]

Admissible sets exist for every n: powers of 3 have 3-progression-free subset sums, as consecutive binary base-3 digit equations show, and hence are k-admissible. The minima are minima of nonempty sets of positive integers.

Corollary 3.2 proves

\[
\begin{adjustbox}{max width=.92\linewidth}
$\displaystyle F_k(m+n)\le F_k(m)F_k(n).$
\end{adjustbox}
\]

Let alpha=inf\_\{m\textgreater=1\} log(F\_k(m))/m. Fix m and write n=tm+r with 0\textless=r\textless m. Repeated composition gives

\[
\begin{adjustbox}{max width=.92\linewidth}
$\displaystyle \frac{\log F_k(n)}n\le
 \frac{t\log F_k(m)+\log F_k(r)}n.$
\end{adjustbox}
\]

The finitely many remainder terms are bounded with m fixed. Hence the limit superior is at most log(F\_k(m))/m, and then at most alpha upon taking the infimum. Every term is at least alpha, so

\[
\begin{adjustbox}{max width=.92\linewidth}
$\displaystyle \lim_{n\to\infty}F_k(n)^{1/n}=\inf_{m\ge1}F_k(m)^{1/m}.$
\end{adjustbox}
\]

This also supplies a direct proof of the subadditivity step without importing a limit theorem as a black box.

The cost and maximum are related by the exact inequalities

\[
\begin{adjustbox}{max width=.92\linewidth}
$\displaystyle 2g_k(n)+1\le F_k(n)\le2ng_k(n)+1.$
\end{adjustbox}
\]

For the left inequality use max(A)\textless=S(A) on a minimizer for F. For the right inequality use S(A)\textless=n max(A) on a minimizer for g. Since g\_k(n)\textgreater=1,

\[
\begin{adjustbox}{max width=.92\linewidth}
$\displaystyle \frac{F_k(n)}{2n+1}\le g_k(n)\le\frac{F_k(n)}2.$
\end{adjustbox}
\]

Taking nth roots proves existence of Lambda\_k and

\[
\begin{adjustbox}{max width=.92\linewidth}
$\displaystyle \Lambda_k=\inf_{n\ge1}F_k(n)^{1/n}.$
\end{adjustbox}
\]

For orientation, the elementary one-shift argument gives Lambda\_k\textgreater1. If B and B+ a belong to a k-AP-free final cube, each maximal a-chain in B has length at most k-2, because adding a extends its end by one. Thus

\[
\begin{adjustbox}{max width=.92\linewidth}
$\displaystyle |B\cup(B+a)|\ge\frac{k-1}{k-2}|B|.$
\end{adjustbox}
\]

Iterating and using \textbar H(A)\textbar\textless=n max(A)+1 proves

\[
\begin{adjustbox}{max width=.92\linewidth}
$\displaystyle \Lambda_k\ge\frac{k-1}{k-2}.$
\end{adjustbox}
\]

The sharper polynomial factor from Korsky\textquotesingle s adaptive recurrence is independent of the existence proof. Also Lambda\_\{k+1\}\textless=Lambda\_k follows from the admissibility inclusion.

\section{5. Finite cyclic certificates are cofinal for the exact rate}\label{n02-finite-cyclic-certificates-are-cofinal-for-the-exact-rate}

Let M\_k be the class of pairs (q,B) satisfying:

\begin{itemize}
\tightlist
\item
  q\textgreater=2 is an integer;
\item
  B is a nonempty finite set of distinct positive integers, with S(B)\textless q;
\item
  D=H(B), reduced in Z/qZ, has no ordered k-term progression with nonzero step.
\end{itemize}

Reduction

\[
\begin{adjustbox}{max width=.92\linewidth}
$\displaystyle \rho_q:D\hookrightarrow\mathbb Z/q\mathbb Z$
\end{adjustbox}
\]

is injective because 0\textless=d\textless=S(B)\textless q. It preserves addition modulo q, not unrestricted integer addition. The equations needed for progression lifting are accounted for next.

\textbf{Lemma 5.1 (all-length cyclic lift).} For (q,B) in M\_k, define

\[
\begin{adjustbox}{max width=.92\linewidth}
$\displaystyle A_t=\bigcup_{j=0}^{t-1}q^jB.$
\end{adjustbox}
\]

Then \textbar A\_t\textbar=t\textbar B\textbar, max(A\_t)=q\^{}(t-1)max(B), and A\_t is k-admissible.

\textbf{Proof.} Since 1\textless=b\textless q, equality q\^{}i b=q\^{}j b\textquotesingle{} with i=q, a contradiction. The equal-level case gives b=b\textquotesingle. Each subset of A\_t chooses a subset in every block, so

\[
\begin{adjustbox}{max width=.92\linewidth}
$\displaystyle H(A_t)=L_t(D,q).$
\end{adjustbox}
\]

The map D\^{}t to L\_t sending (d\_0,...,d\_\{t-1\}) to sum\_j d\_j q\^{}j is a bijection, with inverse the actual base-q digits. Suppose y\_i is a k-AP in this language. Modulo q, the digits form a k-AP; their common difference must be zero. All low digits therefore equal one actual digit c. The map

\[
\begin{adjustbox}{max width=.92\linewidth}
$\displaystyle \{y\in L_t:y\equiv c\pmod q\}\longrightarrow L_{t-1},\qquad
 y\longmapsto(y-c)/q$
\end{adjustbox}
\]

is a bijection, with inverse z-\textgreater c+qz. It sends the progression difference to its exact integer quotient by q. Induction on t forces the original progression to be constant. This proof uses no primality assumption. QED.

Deleting excess generators yields Lambda\_k\textless=q\^{}(1/\textbar B\textbar).

\textbf{Lemma 5.2 (universal finite embedding).} Every k-admissible B gives an element

\[
\begin{adjustbox}{max width=.92\linewidth}
$\displaystyle (2S(B)+1,B)\in M_k.$
\end{adjustbox}
\]

\textbf{Proof.} Put q=2S(B)+1. A modular k-AP among the displayed integer values in D has all second differences divisible by q. Each lies in {[}-2S(B),2S(B){]}, so each vanishes as an integer. The displayed values are therefore an integer k-AP and hence constant. Their modular step is zero. QED.

Let beta=inf\_\{(q,B) in M\_k\} q\^{}(1/\textbar B\textbar). Lemma 5.1 gives Lambda\_k\textless=beta. Lemma 5.2 gives beta\textless=F\_k(n)\^{}(1/n) for every n. Section 4 then gives beta\textless=Lambda\_k. This proves Theorem 1.

For each fixed q the literal search domain is finite: B is a subset of {[}1,q-1{]} and sum B\textless q. Define

\[
\begin{adjustbox}{max width=.92\linewidth}
$\displaystyle m_k(q)=\max\{|B|:(q,B)\in M_k\}.$
\end{adjustbox}
\]

Ignore q with no nonempty certificate. Then

\[
\begin{adjustbox}{max width=.92\linewidth}
$\displaystyle \Lambda_k=\inf_q q^{1/m_k(q)}.$
\end{adjustbox}
\]

The finite sequence obtained by restricting q\textless=Q decreases to Lambda\_k. This is a proved upper approximation; the identity alone supplies no computable error bound after a finite Q. Comparisons between candidate rates q\^{}(1/m) and r\^{}(1/s) are made exactly by comparing q\^{}s and r\^{}m.

\section{6. Full carry graph, including nonstandard digits}\label{n02-full-carry-graph-including-nonstandard-digits}

Let D be any finite set of nonnegative integers containing 0; write W=max D. Fix q\textgreater=2 and k\textgreater=3, and put

\[
\begin{adjustbox}{max width=.92\linewidth}
$\displaystyle L=\left\lceil\frac W{q-1}\right\rceil.$
\end{adjustbox}
\]

A state is (c\_1,...,c\_\{k-1\},b), where b is a Boolean recording whether a nonzero digit of the progression difference has occurred, and

\[
\begin{adjustbox}{max width=.92\linewidth}
$\displaystyle -L\le c_i\le L+i.$
\end{adjustbox}
\]

The initial state is (0,...,0,false). The accepting state is (0,...,0,true).

For a in D, e in \{0,...,q-1\}, and b\_1,...,b\_\{k-1\} in D, draw a labeled edge when the integers c\textquotesingle\_i satisfy

\[
\begin{adjustbox}{max width=.92\linewidth}
$\displaystyle \boxed{a+i e+c_i=b_i+q c'_i\quad(1\le i<k).}$
\end{adjustbox}
\]

The target flag is b OR (e!=0). Labels include the actual a,e,b\_i, not merely their residues. When several actual digits have the same residue, all choices of b\_i are retained.

\subsection{6.1 Invariant bounds}\label{n02-invariant-bounds}

The next carry is (a+i e+c\_i-b\_i)/q. For the lower bound its numerator is at least -W-L. Since W\textless=L(q-1), this quotient is at least -L. For the upper bound its numerator is at most W+i(q-1)+L+i=W+L+iq\textless=q(L+i). Thus the stated box is invariant. The graph has at most

\[
\begin{adjustbox}{max width=.92\linewidth}
$\displaystyle V=2\prod_{i=1}^{k-1}(2L+i+1)$
\end{adjustbox}
\]

states.

When D is contained in {[}0,q-1{]}, b\_i is the unique remainder of a+i e+c\_i. Starting from zero, 0\textless=c\_i\textless=i is invariant, so only

\[
\begin{adjustbox}{max width=.92\linewidth}
$\displaystyle V=2\prod_{i=1}^{k-1}(i+1)=2k!$
\end{adjustbox}
\]

states are needed. The reduction of this bound uses uniqueness of ordinary digit remainders; the unrestricted graph retains signed carries.

\subsection{6.2 Path-to-progression map}\label{n02-path-to-progression-map}

For a length-t accepting path with labels a\_j,e\_j,b\_\{i,j\}, define

\[
\begin{adjustbox}{max width=.92\linewidth}
$\displaystyle x_0=\sum_{j=0}^{t-1}a_jq^j,\quad
 \delta=\sum_{j=0}^{t-1}e_jq^j,\quad
 x_i=\sum_{j=0}^{t-1}b_{i,j}q^j.$
\end{adjustbox}
\]

Multiply each edge identity by q\^{}j and sum over j. Interior carry terms telescope; initial and terminal carries are zero. Hence

\[
\begin{adjustbox}{max width=.92\linewidth}
$\displaystyle x_i=x_0+i\delta.$
\end{adjustbox}
\]

The accepting flag gives delta\textgreater0, and all x\_i belong to L\_t(D,q). This is an explicit witness map from labeled accepting paths to progression representations.

\subsection{6.3 Progression-to-path map}\label{n02-progression-to-path-map}

Conversely, take any positive-step integer k-AP in L\_t. Choose its actual D-word representations. Choose ell\textgreater=t so q\^{}ell exceeds its positive difference, and append zero digits to every representation. Expand the difference in its ordinary ell base-q digits e\_j. Starting at zero, the identities in Section 6 define integer carries, because equality of the final numerical values implies divisibility of each low-prefix residual by the next power of q. Explicitly,

\[
\begin{adjustbox}{max width=.92\linewidth}
$\displaystyle c_{i,j}=\frac{\sum_{h<j}(a_h+i e_h-b_{i,h})q^h}{q^j}.$
\end{adjustbox}
\]

The value is integral: subtract the remaining high-position terms from x\_0+i delta-x\_i=0. These carries satisfy the transition equation, stay in the invariant box, and end at zero. A nonzero difference produces the accepting flag. Negative-step progressions are handled by the involution

\[
\begin{adjustbox}{max width=.92\linewidth}
$\displaystyle (x,d)\longmapsto(x+(k-1)d,-d),$
\end{adjustbox}
\]

which reverses the tuple and has itself as inverse.

For ordinary digits, no padding beyond t is needed, and uniqueness of all digit expansions makes the two maps inverse bijections between accepting paths of length t and positive-step k-APs in L\_t. For unrestricted D, the path labels preserve each chosen representation; the map to numerical progressions is surjective. Its fibers are exactly the available representation choices, which are irrelevant to existence but have not been identified with distinct values.

\textbf{Proof of Theorem 2.} A forbidden progression exists at some finite length exactly when the accepting state is reachable. Reachability in the displayed finite graph is decidable. A shortest path never repeats a state, so a failure has a witness of length at most V-1. An accepting path returns the explicit digits and points above. A safe graph is certified by a reachable set containing the initial state, closed under every labeled transition, and excluding the accepting state. QED.

\subsection{6.4 Relation to modular obstructions}\label{n02-relation-to-modular-obstructions}

For ordinary digits, a nonzero-step modular k-AP is exactly an initial edge with e!=0. Its endpoint consists of the carries floor((a+i e)/q). Thus there is a typed bijection

\[
\begin{adjustbox}{max width=.92\linewidth}
$\displaystyle \operatorname{AP}_{k,\mathbb Z/q\mathbb Z}^*(D)
 \simeq\{\text{initial carry edges with }e\ne0\}.$
\end{adjustbox}
\]

A modular certificate excludes all such initial edges and leaves only the zero/false state reachable. The full graph permits initial nonzero-step edges whose carry states never reach the accepting state.

For a concrete strict example, B=\{2,7\}, D=\{0,2,7,9\}, q=11, k=3. Its four modular progressions have (start,step)=(0,9),(2,9),(7,2),(9,2). The complete reachable set is

\[
\begin{adjustbox}{max width=.92\linewidth}
$\displaystyle ((0,0),false),\quad((0,1),true),\quad((1,1),true).$
\end{adjustbox}
\]

The initial state has four constant self-edges and four edges to the other two states. From (0,1), the labels (a,e) are (2,5),(7,6),(9,4), with targets (0,1),(1,1),(1,1). From (1,1), the labels are (0,6),(2,4),(7,5), with targets (0,1),(0,1),(1,1). These are all 14 valid edges. None reaches ((0,0),true). Therefore every base-11 word with these digits is 3-AP-free despite those modular progressions.

Signed carries are also essential for the unrestricted theorem. With D=\{0,2\}, q=2, k=3, the actual progression (0,2,4) is encoded by the two labels

\[
\begin{adjustbox}{max width=.92\linewidth}
$\displaystyle (a,e,b_1,b_2)=(0,0,2,0),\quad(0,1,0,2).$
\end{adjustbox}
\]

The intermediate carries are (-1,0), followed by (0,0). Excluding negative carries would discard this witness.

\section{7. From unrestricted carry certificates back to finite cyclic certificates}\label{n02-from-unrestricted-carry-certificates-back-to-finite-cyclic-certificates}

For distinct positive B and q\textgreater=2, write each actual b uniquely as q\^{}v u with v\textgreater=0 and q not dividing u. This is a bijection between positive integers and the pairs (v,u), with exact inverse (v,u)-\textgreater q\^{}v u. The pair, including v, is retained.

Call (q,B) separated when distinct members of B have distinct residual coordinates u. Then the sets q\^{}jB have no cross-level collisions: equality q\^{}i b=q\^{}j b\textquotesingle{} is equivalent to u=u\textquotesingle{} and i+v=j+v\textquotesingle. Conversely, equal residual coordinates for different b give a cross-level collision at a sufficiently long prefix.

Let C\_k be the finite pairs (q,B) that are separated and whose full graph for D=H(B) excludes acceptance. Section 6 proves that their prefixes A\_t have exactly t\textbar B\textbar{} positive distinct generators and are k-admissible. Their largest generator is q\^{}(t-1)max(B), so

\[
\begin{adjustbox}{max width=.92\linewidth}
$\displaystyle \Lambda_k\le q^{1/|B|}.$
\end{adjustbox}
\]

Every M\_k certificate is in C\_k. There is also a concrete family of maps from C\_k back to M\_k:

\[
\begin{adjustbox}{max width=.92\linewidth}
$\displaystyle (q,B)\longmapsto
 \left(2S(B)\frac{q^t-1}{q-1}+1,\;\bigcup_{j<t}q^jB\right).$
\end{adjustbox}
\]

The integer in the first coordinate is 2S(A\_t)+1. Lemma 5.2 proves modular admissibility. Its cost per generator satisfies

\[
\begin{adjustbox}{max width=.92\linewidth}
$\displaystyle \left(2S(B)\frac{q^t-1}{q-1}+1\right)^{1/(t|B|)}
 \longrightarrow q^{1/|B|}.$
\end{adjustbox}
\]

Thus the unrestricted carry class and the simpler cyclic class have the same infimum of exponential costs:

\[
\begin{adjustbox}{max width=.92\linewidth}
$\displaystyle \boxed{\Lambda_k=\inf_{(q,B)\in C_k}q^{1/|B|}
 =\inf_{(q,B)\in M_k}q^{1/|B|}.}$
\end{adjustbox}
\]

This equality keeps finite-level differences visible: a particular carry certificate may fail the modular test, while its explicitly constructed longer prefixes approach the same rate through valid modular certificates.

\section{8. Applying the method: the six-generator block}\label{n02-applying-the-method-the-six-generator-block}

Let

\[
\begin{adjustbox}{max width=.92\linewidth}
$\displaystyle B_* = \{1,4,5,17,21,22\},\qquad S(B_*)=70.$
\end{adjustbox}
\]

Its actual subset-sum image is the following 38-element set:

\[
\begin{adjustbox}{max width=.92\linewidth}
$\displaystyle \begin{aligned}
 D_* = \{&0,1,4,5,6,9,10,17,18,21,22,23,25,26,27,28,30,31,32,\\
          &38,39,40,42,43,44,45,47,48,49,52,53,60,61,64,65,66,69,70\}.
\end{aligned}$
\end{adjustbox}
\]

The certificate computes this image directly from all subset choices. It then checks every ordered pair (x,d) with 0\textless=x\textless q and 1\textless=d\textless q. There are 9,312 pairs for q=97 and 8,556 pairs for q=93.

An independent chain calculation gives a compact description of the modular checks. For each nonzero step, decompose D\_* into maximal chains under addition of that step. Full cosets are checked separately by verifying that the chain decomposition covers D\_*; hence repeating progressions in a proper subgroup are included. The histogram of the longest chain per step is:

\begin{longtable}[]{@{}
  >{\raggedleft\arraybackslash}p{(\columnwidth - 8\tabcolsep) * \real{0.2000}}
  >{\raggedleft\arraybackslash}p{(\columnwidth - 8\tabcolsep) * \real{0.2000}}
  >{\raggedleft\arraybackslash}p{(\columnwidth - 8\tabcolsep) * \real{0.2000}}
  >{\raggedleft\arraybackslash}p{(\columnwidth - 8\tabcolsep) * \real{0.2000}}
  >{\raggedleft\arraybackslash}p{(\columnwidth - 8\tabcolsep) * \real{0.2000}}@{}}
\toprule\noalign{}
\begin{minipage}[b]{\linewidth}\raggedleft
q
\end{minipage} & \begin{minipage}[b]{\linewidth}\raggedleft
longest chain 2
\end{minipage} & \begin{minipage}[b]{\linewidth}\raggedleft
longest chain 3
\end{minipage} & \begin{minipage}[b]{\linewidth}\raggedleft
longest chain 4
\end{minipage} & \begin{minipage}[b]{\linewidth}\raggedleft
longest chain 5
\end{minipage} \\
\midrule\noalign{}
\endhead
\bottomrule\noalign{}
\endlastfoot
97 & 6 steps & 28 steps & 62 steps & 0 steps \\
93 & 4 steps & 18 steps & 68 steps & 2 steps \\
\end{longtable}

The rows sum to q-1. Therefore D\_* has no nonzero-step 5-AP modulo 97 and no nonzero-step 6-AP modulo 93. For 93 the argument includes all steps of nontrivial gcd with 93, rather than invoking a prime-modulus lemma.

Lemma 5.1 now proves the all-length constructions. Because \(\max(B_*)=22<q\min(B_*)\), their elements are ordered by level and then by the displayed weights. Taking the first n generators gives exactly the bound in Theorem 3. This proves the theorem for every n, not just multiples of six.

The earlier successful blocks found on the same search route remain useful at small finite sizes:

\begin{longtable}[]{@{}rrlrl@{}}
\toprule\noalign{}
k & q & B & number of subset sums & upper rate \\
\midrule\noalign{}
\endhead
\bottomrule\noalign{}
\endlastfoot
4 & 19 & \{1,7,8\} & 7 & 19\^{}(1/3), existing contributor\textquotesingle s control \\
5 & 23 & \{1,3,4,7\} & 12 & 23\^{}(1/4) \\
6 & 47 & \{1,2,6,7,14\} & 23 & 47\^{}(1/5) \\
5 & 97 & B\_* & 38 & 97\^{}(1/6) \\
6 & 93 & B\_* & 38 & 93\^{}(1/6) \\
\end{longtable}

Taking the minimum of the applicable finite upper expressions is valid. Asymptotically, \(97^{1/6}<23^{1/4}\) because \(97^2<23^3\); \(93^{1/6}<47^{1/5}\) follows from direct integer powers.

\subsection{8.1 A correlated vector model explaining the block}\label{n02-a-correlated-vector-model-explaining-the-block}

Use the exact column map \(V:\mathbb Z^6\to\mathbb Z^3\) with columns

\[
\begin{adjustbox}{max width=.92\linewidth}
$\displaystyle e_1,\ e_2,\ e_1+e_2,\ e_3,\ e_2+e_3,\ e_1+e_2+e_3.$
\end{adjustbox}
\]

The scalar map L:Z\^{}3-\textgreater Z, L(x,y,z)=x+4y+17z, satisfies L(V epsilon)=sum\_i b\_i epsilon\_i for every binary epsilon. Thus the square from binary subsets through the vector image or through the six numerical generators commutes exactly.

Every point of V(\{0,1\}\^{}6) has 0\textless=x,z\textless=3 and 0\textless=y\textless=4. A five-term integer-vector AP has zero x- and z-differences, since a nonzero such difference would span at least four. A nonconstant y progression must run from 0 to 4 or conversely. At y=0, only columns e\_1,e\_3 can occur, giving x,z in \{0,1\}. At y=4, all four columns with nonzero y must occur, giving x,z in \{2,3\}. These endpoint fibers are disjoint. Therefore the vector subset-sum image is 5-AP-free.

The scalar map can create additional additive relations; its action on each proposed modular progression is therefore checked by the complete finite certificate above, rather than assuming arbitrary projections preserve freeness. This correlated model motivates the six-generator choice and identifies the endpoint information lost by replacing its image with the full rectangular box.

\subsection{8.2 Explicit failures of smaller bases}\label{n02-explicit-failures-of-smaller-bases}

For this fixed B\_*, every q from 71 through 96 fails at k=5, and every q from 71 through 92 fails at k=6. These are bounded statements about this exact block; the full graph returns integer witnesses. For example, q=93 at k=5 returns

\[
\begin{adjustbox}{max width=.92\linewidth}
$\displaystyle (9,990,1971,2952,3933),\qquad d=981,$
\end{adjustbox}
\]

all in the two-block language. The q=92, k=6 witness is

\[
\begin{adjustbox}{max width=.92\linewidth}
$\displaystyle (1,967,1933,2899,3865,4831),\qquad d=966.$
\end{adjustbox}
\]

The k=5 base-22 version of \{1,3,4,7\} fails with (0,77,154,231,308). The published receipt retains the actual digit labels for all 57 nearby-base obstruction cases, including the familiar k=4 failures at bases 17 and 18.

\section{9. A proved bound on the rectangular-digit route}\label{n02-a-proved-bound-on-the-rectangular-digit-route}

Let c\_1,...,c\_m be distinct nonzero vectors in N\^{}r with coordinate totals sum\_j c\_\{j,i\}\textless=tau. At most r columns have l1 mass one. Every other column has mass at least two. Consequently

\[
\begin{adjustbox}{max width=.92\linewidth}
$\displaystyle 2m-r\le\sum_j\|c_j\|_1\le r\tau,
 \qquad\boxed{m\le\lfloor r(\tau+1)/2\rfloor.}$
\end{adjustbox}
\]

This is the count attained by the singleton-plus-edge construction in Korsky\textquotesingle s appropriate nearly-regular graph range. For p\textgreater tau the evaluation

\[
\begin{adjustbox}{max width=.92\linewidth}
$\displaystyle E_p:\{0,\ldots,\tau\}^r\hookrightarrow\mathbb N,\qquad
 x\longmapsto\sum_i x_i p^i$
\end{adjustbox}
\]

is injective, with inverse its displayed digit vector. It sends the vector subset-sum image to the integer subset-sum image and retains the generator count. Its full-box envelope discards correlations between coordinates. Section 8.1 exhibits the precise lost data: the fibers at y=0 and y=4 have disjoint outer coordinates even though the full box would allow a five-term vertical progression.

Thus the column-mass bound is a proved obstruction to improving generator density while retaining that particular rectangular resource constraint. The carry graph and the correlated block operate on the actual image and keep the discarded endpoint constraints.

\section{10. Exact remaining general question and research implications}\label{n02-exact-remaining-general-question-and-research-implications}

Korsky\textquotesingle s inspected bound gives

\[
\begin{adjustbox}{max width=.92\linewidth}
$\displaystyle \frac{k-1}{k-2}\le\Lambda_k\le
 \min_{p\ge3\text{ prime}}p^{2/(\min\{p,k\}-1)}.$
\end{adjustbox}
\]

Theorem 1 replaces the separate lower/upper limits with an existing limit. Theorem 3 tightens the right side at k=5 and k=6. For large k, the inspected bounds still leave logarithmic rates between order 1/k and order log(k)/k.

In the new exact finite formulation, the desired lower scale is the uniform statement

\[
\begin{adjustbox}{max width=.92\linewidth}
$\displaystyle \frac{\log q}{|B|}\ge c\frac{\log k}{k}
 \quad\text{for every }(q,B)\in M_k$
\end{adjustbox}
\]

for some absolute positive c in the stated large-k regime. This is a research target, not a premise used in any proved claim here. A violation comes with a finite pair (q,B) and, by Section 5, an explicit infinite family of integer constructions. A proof must control every finite cyclic certificate; checking finitely many q cannot establish that quantifier. The exact maps in Sections 3, 5, 6, and 7 make both outcomes usable without discarding carries or changing the generator problem.

The next mathematical objective is therefore a uniform inverse or packing inequality for these actual subset-sum images. The rectangular count in Section 9 settles one restricted representation class; the six-generator examples demonstrate why its envelope alone cannot supply the general answer. The certificate search remains a productive way to test proposed lower inequalities and return concrete counterexamples when they fail.

\section{11. Verification and reproducibility}\label{n02-verification-and-reproducibility}

The standard-library checker is \nolinkurl{certificates/verify_general_k_capacity.py}. Run

\begin{Verbatim}[breaklines=true,breakanywhere=true,fontsize=\footnotesize]
python certificates/verify_general_k_capacity.py \
  --output certificates/general_k_capacity_receipt.json
\end{Verbatim}

The default run does all of the following:

\begin{itemize}
\tightlist
\item
  audits the five displayed finite blocks by direct modular enumeration, independent maximal-chain calculations, full carry reachability, and direct actual generator subsets through two blocks;
\item
  checks every D subset {[}0,q-1{]} containing 0 for q=2,...,6 and k=3,...,6: 248 alphabets, 128 safe and 120 unsafe, with 744 direct language-length comparisons;
\item
  checks every D subset {[}0,2q{]} containing 0 and max(D)\textgreater=q, for q=2,3,4 and k=3,4,5: 966 alphabets, 29 safe and 937 unsafe, with 1,932 direct comparisons;
\item
  checks 738 literal compositions from 246 admissible left instances in the stated finite domain, preserving the exact sumset bijection, source cardinality, cost product, and progression exclusion;
\item
  replays all 57 nearby-base failure witnesses, the prime-base carry-only example, signed carry failure, nonstandard-digit safe example, and cross-level generator-collision detection;
\item
  performs a literal branch-and-bound search for modular blocks for every q=3,...,48 and k=5,6, with no rescaling of candidate weights.
\end{itemize}

Every safe graph is rechecked for closure under all valid labels. Failure witnesses are reconstructed by integer sums and checked against the actual language. No floating-point arithmetic is used in a proof check. Search upper-rate comparisons use integer powers. A transcript hash fixes each deterministic bounded replay.

Exploratory C++ scans additionally supplied the candidates at q=93 and 97. Candidate discovery is separated from verification: the reproducible Python certificate checks those blocks directly without relying on the exploratory search or on a claim of optimality at those moduli. The command-line option \texttt{-\/-search-max-base\ 127} permits a larger fully replayed search; the default receipt only claims the range it actually executes.

The receipt records its precise source hash. The companion summary records that hash and the full receipt hash, together with the principal results and chain profiles. The downloadable contribution bundle includes the full receipt. Existing workbench validation and Lean receipts have their earlier scopes; they are not silently extended to these files.

\subsection{Formalization units}\label{n02-formalization-units}

The direct finite modules are: composition and its inverse; second-difference decomposition; the cost product and infimum limit proof; modular lifting with composite q; unrestricted carry-state invariance; accepting-path/represented-progression correspondence; separation of geometric generators; finite reachable-set closure; and the two six-generator certificates. All new units remain outside current Lean coverage. They can be formalized independently of PR \#2\textquotesingle s signed-block variance development.

\section{References and source roles}\label{n02-references-and-source-roles}

Anonymous/deleted Reddit contributor, \& The Clankers. (2026, September 13). \emph{The exponential rate for four-term-progression-free subset-sum sets} {[}Workbench manuscript{]}. \texttt{paper/ep817\_k4\_rate.tex} at the base revision. Role: existing k=4 construction and algebraic comparison; original mathematical attribution retained.

Korsky, S. (2026, June 23). \emph{Arithmetic progression-free subset-sum sets} (arXiv:2606.24139v1). \url{https://arxiv.org/html/2606.24139v1} . Inspected: Theorems 1.2, 1.3 and 6.3; Sections 5--7. Role: benchmark general bounds, exact conventions, one-shift context, rectangular graph construction, and large-k target.

Costa, S. (2026, September 5). \emph{A negative answer to the Erdős--Sárkőzy question} (arXiv:2609.06303v1). \url{https://arxiv.org/html/2609.06303v1} . Inspected: Theorem 1.1 and the reduction. Role: separates the historical k=3 positive-prefactor question from the broader all-k rate problem. Its construction is not an input to the proofs above.

sneed-and-feed. (2026, September 13). \emph{Formalize finite upper bound theorem and core helper lemmas in Extended.lean} {[}Pull request \#1{]}. \url{https://github.com/KokunoYumeto/erdos-problem-817-workbench/pull/1} . Role: existing formal upper-bound contribution, with credit separate from the anonymous mathematical construction.

The Clankers. (2026, September 9). \emph{A handoff describes the research; it does not command the next researcher} {[}PolyClank design document{]}. \nolinkurl{docs/polyclank/RESEARCH_STATE.md} in \nolinkurl{KokunoYumeto/yang-mills-interacting-workbench} at \nolinkurl{7ba7ea5a5dab84d88934477de93dda71f6af0662}. Role: inspected research-record convention; no mathematical dependency or peer-validation claim.

Carpenter, R., Defant, C., \& Kravitz, N. (2026, May 6). \emph{Sets with few subset sums} (arXiv:2605.05498v1). \url{https://arxiv.org/html/2605.05498v1} . Abstract and source entry inspected as a possible inverse-theorem direction; not used in a proof here.

Croot, E., Mao, J., \& Yip, C. H. (2026, March 15). \emph{Hilbert cubes in sets with arithmetic properties} (arXiv:2603.14654v1). \url{https://arxiv.org/html/2603.14654v1} . Introduction/framework statements inspected; no theorem imported into the present proof.
